\documentclass[a4paper,11pt]{article}
\usepackage{exam-style-842-unit}
\examinehead{842 信息化概论 · 第 05 单元练习册}{练习一\quad 云计算基础}
\begin{document}

\unitcover{842 信息化概论 · 单元练习册}{第 05 单元\quad 云计算与大数据}{练习一\quad 云计算基础（共 100 分）}{约 75 分钟}{云计算的由来、定义与特征；三种服务模式；三种部署模式与社区云；云计算的核心技术；云计算与网格计算、分布式计算的联系与区别。}

\begin{knowmap}
\textbf{由来：}20 世纪 60 年代麦卡锡提出“计算能力像水和电一样作为公用事业”$\rightarrow$ 1999 年 Salesforce.com 通过网站向企业提供企业级应用 $\rightarrow$ 2002 年亚马逊提供 Web Service、2005 年提出弹性计算云 EC2 $\rightarrow$ 2006 年 8 月 Google 首席执行官施密特在搜索引擎大会上首次提出“云计算”\\[2pt]
\textbf{定义：}教材称云计算是基于互联网的相关服务的增加、使用和交付模式，通过互联网提供动态易扩展且经常是虚拟化的资源；NIST 称云计算是按使用量付费的模式，提供按需网络访问和可配置的计算资源共享池\\[2pt]
\textbf{特征：}教材 4.1.3 列七个特点——超大规模、虚拟化、高可靠性、通用性、高可扩展性、按需服务、廉价；NIST 另概括五个基本特征——按需自助服务、广泛网络接入、资源池化、快速弹性、可计量服务\\[2pt]
\textbf{服务模式：}按 NIST 划分为 IaaS（第 1 层）$\rightarrow$ PaaS（第 2 层）$\rightarrow$ SaaS（第 3 层）三大类\\[2pt]
\textbf{部署模式：}公有云 · 私有云 · 混合云（教材列 3 种）；NIST 另有社区云，面向有共同关切的特定社区，教材未展开\\[2pt]
\textbf{核心技术：}虚拟化技术、大规模数据处理技术、云计算平台管理技术、绿色节能技术；云计算是网格计算、分布式计算、并行计算、效用计算、网络存储、虚拟化、负载均衡等技术融合的产物
\end{knowmap}

\begin{keyterms}
\textbf{云计算}：基于互联网的相关服务的增加、使用和交付模式，通过互联网提供动态易扩展且经常是虚拟化的资源。\\
\textbf{NIST 云计算定义}：一种按使用量付费的模式，提供可用的、便捷的、按需的网络访问，进入可配置的计算资源共享池。\\
\textbf{虚拟化}：在软件中仿真计算机硬件、以虚拟资源为用户提供服务的计算形式，旨在合理调配资源、提高利用效率。\\
\textbf{按需服务}：云是一个庞大的资源池，可按需购买，并像自来水、电、煤气那样计费。\\
\textbf{IaaS}：基础设施即服务，位于云计算最底层，向用户提供服务器、存储、网络等基本计算资源。\\
\textbf{PaaS}：平台即服务，位于第二层，向用户提供分布式操作系统、分布式数据库等基础软件。\\
\textbf{SaaS}：软件即服务，位于第三层，用户通过网页浏览器或终端使用远程服务器上的应用。\\
\textbf{公有云}：云服务提供商为公众提供的云计算平台，理论上任何人都可通过授权接入。\\
\textbf{私有云}：云服务提供商为企业在其内部建设的专有云计算系统，存在于防火墙内，只为企业内部服务。\\
\textbf{混合云}：同时提供公有和私有服务的云计算系统，是介于公有云和私有云之间的折中方案。\\
\textbf{社区云}：NIST 定义的第 4 种部署模式，为具有共同关切（如安全、合规）的特定社区提供云服务，教材未列入。\\
\textbf{云安全}：通过大量客户端监测网络中软件行为的异常，把分析结果和病毒木马解决方案分发到各客户端。\\
\textbf{绿色节能技术}：通过规模经济效益和节能措施降低数据中心能耗，是云计算核心技术之一。
\end{keyterms}

\begin{contrast}{云计算 / 网格计算 / 分布式计算（每行依次为：核心思想·资源组织方式·典型目标；后两者教材未给定义，按通行表述整理）}
\citem{云计算}{通过互联网按需提供可配置的计算资源，并按使用量付费}{资源池化、弹性伸缩、多用户共享}{让计算像水电一样取用方便、成本低廉}
\citem{网格计算}{把地理上分布、异构的计算资源集成为一台“虚拟超级计算机”}{跨机构共享空闲资源、统一调度}{解决大规模科学计算与协同研究问题}
\citem{分布式计算}{把计算任务分解为若干子任务，分配到多台计算机上并行处理}{多节点协同、任务分解与结果汇总}{提高计算速度与系统吞吐量}
\end{contrast}

\begin{contrast}{公有云 / 私有云 / 混合云（每行依次为：服务对象·建设与部署位置·优缺点）}
\citem{公有云}{面向外部客户，理论上任何人授权即可接入}{互联网上发布，由云服务商建设运营}{发挥规模经济效益，但增加安全风险，存在法规、计费、迁移等缺陷}
\citem{私有云}{为“一个”客户单独使用而构建}{部署在企业数据中心的防火墙内或安全的主机托管场所}{对数据、安全性和服务质量控制最有效，但成本高、规模经济效益受限}
\citem{混合云}{同时面向内部与外部客户}{公有服务与私有服务并存}{兼顾安全与弹性，是介于公有云与私有云之间的折中方案}
\end{contrast}

\begin{contrast}{虚拟化的两种应用模式（教材 4.1.5；每行依次为：做法·目的·典型场景）}
\citem{一虚多}{把一台性能强大的服务器虚拟成多个独立的小服务器}{服务不同用户，提高资源利用率}{服务器整合、多用户共享同一物理设备}
\citem{多虚一}{把多台服务器虚拟成一台强大的服务器}{集中算力完成特定功能}{高性能计算、大规模集群任务}
\end{contrast}

\qtypename{一、单项选择题}{10 题，每题 3 分，共 30 分}
\q 早在 20 世纪 60 年代就提出把计算能力作为一种像水和电一样的公用事业提供给用户的是（\quad）
\choicesline{A. 麦卡锡}{B. 施密特}{C. 比尔·盖茨}{D. 摩尔}

\q 2006 年 8 月在搜索引擎大会上首次提出“云计算”概念的是（\quad）
\choicesline{A. 施密特}{B. 贝索斯}{C. 卡廷}{D. 扎克伯格}

\q 按美国国家标准与技术研究院（NIST）的定义，云计算是一种按（\quad）付费的模式。
\choicesline{A. 使用量}{B. 服务器数量}{C. 用户人数}{D. 软件套数}

\q 用户无需了解应用运行的具体位置，只需一台计算机或手机即可通过终端获取服务，这体现了云计算的（\quad）特点。
\choicesline{A. 超大规模}{B. 虚拟化}{C. 高可靠性}{D. 通用性}

\q “云”的规模可以动态伸缩，满足应用和用户规模增长的需要，这体现了云计算的（\quad）特点。
\choicesline{A. 通用性}{B. 高可靠性}{C. 高可扩展性}{D. 廉价}

\q 下列不属于 NIST 概括的云计算五个基本特征的是（\quad）
\choicesline{A. 按需自助服务}{B. 广泛网络接入}{C. 快速弹性}{D. 数据独占性}

\q 企业内部（专网）发布、只为企业内部服务的云计算系统是（\quad）
\choicesline{A. 公有云}{B. 私有云}{C. 混合云}{D. 社区云}

\q 按 NIST 划分，云计算的服务模式由低到高依次是（\quad）
\choicesline{A. SaaS、PaaS、IaaS}{B. PaaS、IaaS、SaaS}{C. IaaS、PaaS、SaaS}{D. IaaS、SaaS、PaaS}

\q 教材指出公有云存在若干“致命缺陷”，下列不属于其缺陷的是（\quad）
\choices{A. 政府没有运营云的法律法规，数据出问题难以追究责任}{B. 云运营商计费标准不统一，流量计费千差万别}
{C. 能给集团企业提供云服务的运营商数量还不多}{D. 更换云服务商时数据迁移比私有云更容易}

\q 教材指出，为云计算服务提供基础架构层面支撑、被称为云计算最重要核心技术之一的是（\quad）
\choicesline{A. 虚拟化技术}{B. 绿色节能技术}{C. 分布式存储技术}{D. 负载均衡技术}

\qtypename{二、多项选择题}{4 题，每题 4 分，共 16 分；每题有两个或两个以上正确项}
\q 教材 4.1.3 归纳的云计算特点包括（\quad）
\choicesline{A. 超大规模}{B. 高可靠性}{C. 通用性}{D. 数据独占性}

\q 关于公有云与私有云，下列说法正确的有（\quad）
\choicesline{A. 公有云理论上任何人都可通过授权接入}{B. 私有云存在于防火墙之外，可为多个外部客户共享}
{C. 私有云安全性更好但成本更高、规模经济效益受限}{D. 混合云是同时提供公有和私有服务的折中方案}

\q 云计算的关键核心技术包括（\quad）
\choicesline{A. 虚拟化技术}{B. 大规模数据处理技术}{C. 云计算平台管理技术}{D. 磁盘碎片整理技术}

\q 云计算是下列哪些传统技术发展融合的产物（\quad）
\choicesline{A. 网格计算}{B. 并行计算}{C. 效用计算}{D. 生物计算}

\qtypename{三、判断题}{6 题，每题 2 分，共 12 分；正确的画“√”，错误的画“×”并写出改正后的表述}
\q 2006 年 8 月，Google 首席执行官施密特在搜索引擎大会上首次提出“云计算”概念。\quad（\quad）\quad 改正：\fixblank

\q 云计算就是虚拟化。\quad（\quad）\quad 改正：\fixblank

\q 公有云与私有云的区别可从建设地点、协议开发程度、服务对象三方面划分。\quad（\quad）\quad 改正：\fixblank

\q 按 NIST 划分，云计算的服务模式有 IaaS、PaaS、SaaS 和 DaaS 四大类。\quad（\quad）\quad 改正：\fixblank

\q 私有云存在于企业防火墙之内，只为企业内部服务，安全性更好但成本更高。\quad（\quad）\quad 改正：\fixblank

\q 教材认为云计算是一种全新创造的共享基础架构方法，与传统计算机技术无关。\quad（\quad）\quad 改正：\fixblank

\qtypename{四、配对题}{6 对，每对 1 分，共 6 分；把 A 栏与 B 栏对应内容配对}
\q 将下列人物、事件或概念与其对应内容配对，把 B 栏字母填入括号。
\par\vspace{0.2em}
\noindent\begin{tabular}{@{}p{0.44\textwidth}@{\hspace{1.2em}}p{0.44\textwidth}@{}}
（1）麦卡锡 \hfill（\quad） & A. 2006 年 8 月在搜索引擎大会上首次提出“云计算”概念 \\
（2）施密特 \hfill（\quad） & B. 20 世纪 60 年代提出计算能力应像水和电一样作为公用事业 \\
（3）Salesforce.com \hfill（\quad） & C. 1999 年提出通过网站向企业提供企业级应用 \\
（4）亚马逊 \hfill（\quad） & D. 2002 年提供 Web Service，2005 年提出弹性计算云 EC2 \\
（5）IaaS \hfill（\quad） & E. 位于第三层，用户通过网页浏览器使用远程服务器上的应用 \\
（6）SaaS \hfill（\quad） & F. 位于最底层，向用户提供服务器、存储、网络等基础设施服务 \\
\end{tabular}\par\vspace{0.3em}

\qtypename{五、简答题}{3 题，每题 8 分，共 24 分}
\q 什么是云计算？写出教材的定义与 NIST 的定义，并简述教材 4.1.3 归纳的七个特点。\anslines{2.4cm}

\q 云计算的核心技术有哪几项？简述各项技术的主要内容与作用。\anslines{2.4cm}

\q 云计算与网格计算、分布式计算有何联系与区别？\anslines{2.4cm}

\newpage
\qtypename{六、综合题}{1 题，共 12 分}
\q 某农业科研院所拟建设科研数据云平台：院内育种、土壤、气象数据涉密，仅供院内使用；同时面向社会公众提供农技问答和农产品行情查询服务；平台还要支撑多个课题组的应用开发。请回答：
（1）该平台应选择哪种（些）部署模式？说明理由。（4 分）
（2）平台的底层计算资源、开发运行环境、面向公众的问答服务分别对应哪种服务模式？（4 分）
（3）有人说“买了虚拟化软件就是建成了云计算”，请结合云计算的核心技术及其与网格计算、分布式计算的关系加以辨析。（4 分）\anslines{5cm}

\end{document}
